% !TeX program = xelatex
% 文档：SmarTAI 今日工作日志（承接《题目识别与判卷问题记录.tex》）
\documentclass[12pt,a4paper]{ctexart}

\usepackage[margin=2.5cm]{geometry}
\usepackage{booktabs}
\usepackage{enumitem}
\usepackage{amsmath}

\title{\bfseries SmarTAI 今日工作日志}
\author{来源：\texttt{SmarTAI-shared} 目录}
\date{\today}

\begin{document}
\maketitle
\sloppy

\section*{今日概览}

\begin{center}
\begin{tabular}{@{}p{3.2cm}p{11.2cm}@{}}
\toprule
类别 & 条目 \\
\midrule
今日成果 & ① 两问题目拆成两道　② 满分分配混乱　③ 选择题不应有步骤分 \\
注意事项 & ① 题目只按题号致各题型“第 1 题”拥挤　② 批改速度慢（3 人卷约 1h）　③ 复核建议加“显示答案” \\
新问题   & ① 会被人工答案误导　② 模型判卷能力不均致低置信度　③ 批改页“运行中”恒为 0 \\
\bottomrule
\end{tabular}
\end{center}

\section*{今日成果}

今日针对昨日《题目识别与判卷问题记录》中列出的问题完成如下修复。

\section{解决昨日问题 1：两问题目拆成两道}

\subsection{问题背景（昨日记录第 2 节）}
同一道解答题（含两问）被重复识别为“整题”与“第一问”两个条目，或一个题目条目
内含两道独立的题，导致同一题被识别为两个、题目总数虚高、判卷重复。

\subsection{今日处理}
\begin{itemize}[leftmargin=2em]
  \item “两问题目”统一拆成两道独立的题目（拆分后各有独立的题号与满分）；
  \item 拆题后对整个题目集做去重：同一题号、同一题面只保留唯一条目；
  \item 若题目本身是整题带子问 (1)(2)(3)，则整题作为一个题目，子问作为内部结构，
        不再单独成条目（与昨日记录的期望一致）；
  \item 代码佐证：\texttt{backend/tools/problem\_dedup.py} 的
        \texttt{dedupe\_extracted\_problems()} 在分数策略冻结前对抽取行做去重
        （含“第一问被拆成独立题”场景），并配套回归测试
        \texttt{backend/tests/test\_problem\_dedup.py}（22 条原始行 → 19 道题）。
\end{itemize}

\section{解决昨日问题 2：由问题 1 引发的满分分配混乱}

\subsection{问题背景（昨日记录第 3 节）}
重复条目各自带一份满分，导致单题满分与卷面总分不匹配、同一题被重复计分。

\subsection{今日处理}
\begin{itemize}[leftmargin=2em]
  \item 每道题只存在一个 \texttt{max\_score}，由教师统一配置；
  \item 识别阶段不再为重复条目擅自生成满分；拆题去重后满分只落在唯一题目上；
  \item 总分统计基于去重后的题目集合，不再重复计入；
  \item 代码佐证：去重发生在 \texttt{resolve\_question\_score\_policy()} 冻结满分
        \textbf{之前}（\texttt{backend/tools/problem\_dedup.py} 头注释），
        保证每题唯一满分（测试 \texttt{test\_deduped\_extraction\_feeds\_score\_policy}）。
\end{itemize}

\section{解决昨日问题 3：选择题不应有步骤分}

\subsection{问题背景（昨日记录第 1 节）}
选择、填空类题目被误判出“过程分”。

\subsection{今日处理}
\begin{itemize}[leftmargin=2em]
  \item 选择题、填空题只判结果分：正确得满分，错误得 0 分，不再输出过程分项；
  \item 仅当题干明确要求写出简要过程时，才在题目类型与评分规则中显式启用过程分；
  \item 代码佐证：\texttt{backend/skills/objective.py} 的
        \texttt{ObjectiveSkill.grade} 明确“result-based, no process scores”，
        返回 \texttt{steps=[]}（客观题永不携带过程分）。
\end{itemize}

\section*{注意事项}

\begin{enumerate}[leftmargin=2em,label=\arabic*.]
  \item \textbf{环节 3（审核题目）与环节 5（校对作答）中，题目只按题号排列}，
        导致选择题的“第 1 题”、填空题的“第 1 题”、解答题的“第 1 题”挤在一起，
        三者共用同一编号难以区分（题号在各题型内重置：单选 1--8、多选 1--3、
        填空 1--3、解答 1--5）。
        代码佐证：各列表均按 \texttt{number} 数值排序
        （如 \texttt{frontend/app/src/routes/tasks/QuestionPreparationOverviewPage.tsx}、
        \texttt{GradingPreflightPage.tsx} 的 \texttt{compareProblems}、
        \texttt{ReviewDetailPage.tsx} 的题目导航）。
        建议：按“题型 + 题号”分组编号，或在题干前显式标注题型，便于区分。

  \item \textbf{环节 6（执行批改）速度仍然较慢}：3 人份卷子已需约 1 小时。
        后续可从并发扇出、调用复用、单次调用超时、失败重试/退避等方面继续优化
        （相关逻辑见 \texttt{backend/agents/multi\_expert.py}、
        \texttt{backend/agents/grading\_agent.py}）。

  \item \textbf{环节 7（复核批改）建议增加“显示答案”功能}：
        人工复核时能对照参考答案，提高复核效率与准确率。
        代码佐证：\texttt{frontend/app/src/routes/tasks/ReviewDetailPage.tsx} 的复核卡片
        目前只展示题干、学生作答、评分标准与批改结果，未渲染
        \texttt{reference\_answer}。
\end{enumerate}

\section*{新问题}

\section{会被人工答案误导}

\subsection{现象}
判卷结果可能被人工作答/人工答案误导，从而对学生的试卷给出错误评分。

\subsection{解决方案}
\begin{enumerate}[leftmargin=2em,label=(\arabic*)]
  \item 增加 AI 复核环节：由 AI 对批改结果进行二次核对；
  \item 在注意事项中明确提示，避免人工答案干扰自动批改。
\end{enumerate}

\section{不同模型判卷能力不均导致低置信度}

\subsection{现象}
不同模型能力不均衡：千问（Qwen）对部分解答题判不了，
DeepSeek（DS）只能判大题。当某题只有部分模型能判时，出现低置信度。

\subsection{解决方案}
改变“低置信度”的合并机制——只要有一个模型给出高置信度，
即视为高置信度结果，不再要求所有模型一致后才判定为高置信度。
（现状：\texttt{backend/agents/multi\_expert.py} 中多专家合并采用
置信度加权平均，且仅剩单个成功样本时降级为
\texttt{degraded\_to\_single} 并强制人工复核，见
\texttt{grading\_agent.py} 的 \texttt{\_apply\_low\_confidence\_policy}。）

\section{执行批改页面中的“运行中”始终为 0}

\subsection{现象}
环节 6（执行批改）页面中的“运行中”计数一直显示为 0，
即便批改任务实际正在执行中也不同步更新。

\subsection{初步定位（代码查证）}
\begin{enumerate}[leftmargin=2em,label=(\arabic*)]
  \item 前端 \texttt{frontend/app/src/routes/tasks/GradingProgressPage.tsx}
        的 \texttt{deriveQueue} 用 \texttt{progress.active} 计算“运行中”数量：
        \texttt{running = min(total - completed, activePairs.size)}，
        因此只要 \texttt{active} 为空，运行的显示恒为 0；
  \item 后端 \texttt{backend/services/task\_facade.py} 的
        \texttt{\_grading\_progress()} 在任务处于 \texttt{grading} 状态时
        返回进度对象，其中 \texttt{"active": []} 是硬编码的空列表；
  \item \texttt{async\_task\_state()} 会把内存 reporter 的真实快照
        （含 \texttt{active}）用上述持久化投影覆盖掉，导致接口返回的
        \texttt{progress.active} 恒为空 → 前端“运行中”恒为 0。
\end{enumerate}

\subsection{解决方案（建议）}
\begin{enumerate}[leftmargin=2em,label=(\arabic*)]
  \item 在 \texttt{\_grading\_progress()} 的持久化投影中加入真实的 active 单位
        （或至少回填 \texttt{active\_units} 计数），而不是硬编码空列表；
  \item 复核 \texttt{async\_task\_state()} 中内存快照被覆盖的逻辑，
        避免“运行中”实时信息在接口层丢失；
  \item 修复后回归验证：批改进行中应显示实际运行题数/人数，完成后归零。
\end{enumerate}

\section*{明日计划（建议）}
\begin{itemize}[leftmargin=2em]
  \item 按“题型 + 题号”重做环节 3/5 的题目编号显示；
  \item 继续优化环节 6 批改速度（并发与重试策略）；
  \item 为环节 7 复核增加“显示答案”功能；
  \item 落实低置信度合并新机制并回归验证；
  \item 修复执行批改页面“运行中”始终为 0 的问题。
  \item 解决解答审核“按下葫芦浮起瓢”问题。
\end{itemize}

\end{document}
